% 36_body.tex — 산출물 ㊱ 엔터프라이즈 딜 데스크 시트 (빈 템플릿 / 완성 예시 공용 본문) — gen_product_templates.py 가 Product_specs.py 에서 생성
% 드라이버: 36_딜데스크시트_빈템플릿.tex (\wbblanktrue) / 36_딜데스크시트_완성예시.tex (\wbblankfalse — 워크북 §X.5 확정 후)
\documentclass[10pt,a4paper]{article}
\usepackage[a4paper,margin=16mm,top=14mm,bottom=20mm]{geometry}
\def\DocNum{36}\def\DocDay{8}\def\DocIdx{1}
\def\DocTitle{엔터프라이즈 딜 데스크 시트}
\def\DocSub{Enterprise Deal Desk — 47 Requests, 5 Dimensions, 3 Classes}
\def\DocVariant{\B{빈 템플릿}{딥게이지 완성 예시}}
\usepackage{wbstyle}
\def\DocCirc{㊱}   % newunicodechar(㉛~㊵ 폴백) 뒤에 정의해야 활성 문자로 읽힌다
\begin{document}
\docheader
 {\Bg{[회사명 · 고객사 2법인]}{딥게이지 · ㈜온담F\&B홀딩스 / ㈜온담푸드시스템}}
 {\Bg{[v1.x / 2028-05]}{v1.0 / 2028-05-10 (D+800)}}
 {\Bg{[작성자 — CPO·CEO]}{강민수 CEO · 윤하경 CPO\rd{7mm}}}
 {\Bg{[승인자 — 전원]}{강민수 · 오세진 · 윤하경\rd{7mm}}}

%% ------------------------------------------------------------------ 1
\secbar{1. 딜 개요}
\guide{온담 21억/3년(연 7억) · 2법인 · RFI 2027-11-02 · 47건/138 mm · 218문항 · 온프렘 · BATNA 대한기공. 정본 §5, E-01~17.}
{\footnotesize
\begin{tabularx}{\linewidth}{|L{36mm}|Y|L{24mm}|}\hline
\hd{항목} & \hd{내용} & \hd{정본 ID} \\ \hline
\ifwbblank
\rd{8mm}계약 총액 / 기간 / 청구 &  &  \\ \hline
\rd{8mm}계약 주체 &  &  \\ \hline
\rd{8mm}라이선스(실질 ARR 증분) &  &  \\ \hline
\rd{8mm}구축·이관 / 온프렘 운영 &  &  \\ \hline
\rd{8mm}요구 47건 — 저자별 건수 &  &  \\ \hline
\rd{8mm}보안·SLA·배상 요구 &  &  \\ \hline
\rd{8mm}BATNA &  &  \\ \hline
\else
계약 총액 / 기간 / 청구 & 21.00억 / 3년 = 구축·이관 3.00(1년차 일시) + 라이선스 11.34 + 온프렘 운영 6.66. 청구 연 7.00억 균등 & E-01·02, §5.2 \\ \hline
계약 주체 & ㈜온담F\&B홀딩스(이천 12 + 베이커리 12 + 61점 30 = 54라인) / ㈜온담푸드시스템(조성태 — 안성 16라인). 두 법인, 두 서명 & §5.1·5.3 \\ \hline
라이선스(실질 ARR 증분) & 연 3.78억 = 라인 환산 70 × 45만 × 12. 70 = 이천 12 / 안성 16 / 베이커리 12 / 61점 30(2개점 = 1라인 — 정책 밖 임시 합의) & E-04, §5.2 \\ \hline
구축·이관 / 온프렘 운영 & 구축 3.00(커버리지 95\% 가정) → 실소요 검산 6.12(실측 61.4\%, D+690) / 운영 연 2.22 = 월 1,850만(GPU 상각 400 + 0.7 FTE 467 + 대기·검증 333 + 회선 150 + 예비 500) & E-03·05·14 \\ \hline
요구 47건 — 저자별 건수 & 나영선 19(62 mm) · 한동석 11(31) · 박세연 9(28) · 오정근 5(12) · 최영주 3(5) = 47건 / 138 mm. 온담 공문은 47건을 '21억에 포함'으로 표기 & §4.12, E-10 \\ \hline
보안·SLA·배상 요구 & 218문항(ISMS-P 미인증, 5영업일 — 회신 D+617: 131 / 62 / 25) · SLA 99.5\% · 배상 온담 100\% vs 초안 10\% · GPU 2대 + 0.7 FTE · 릴리스 분기 1회 승인 후 · CFO 결재 3주 & E-12·13·14·16 \\ \hline
BATNA & [대한기공] 표준계약 연 2.4억 × 3년 = 7.2억(21억의 34.3\%) · 커스텀 0 · 클라우드 · 4주 · [H사] 김도윤 소개 & E-15, §5.8 \\ \hline
\fi
\end{tabularx}}

%% ------------------------------------------------------------------ 2
\landscapeon
\secbar{2. 코드체계 불일치 5차원 표 — 47건을 판정하기 전에 먼저}
\guide{판정 라벨 / 검사항목 식별자 / 시간·수량 단위 / 규격 근거 / 부적합 흐름. 각 차원의 불일치 성격과 1차 판정(제품화 후보·유상 커스텀·거절)과 근거.}
{\footnotesize
\begin{tabularx}{\linewidth}{|L{30mm}|Y|Y|Y|Y|}\hline
\hd{차원} & \hd{자동차 3세트} & \hd{온담 4세트} & \hd{불일치의 성격} & \hd{1차 판정·근거} \\ \hline
\ifwbblank
\rd{13mm}① 판정 라벨 &  &  &  &  \\ \hline
\rd{13mm}② 검사항목 식별자 &  &  &  &  \\ \hline
\rd{13mm}③ 시간·수량 단위 &  &  &  &  \\ \hline
\rd{13mm}④ 규격 근거 &  &  &  &  \\ \hline
\rd{13mm}⑤ 부적합 처리 흐름 &  &  &  &  \\ \hline
\else
① 판정 라벨 & OK / NG 2값 & in/out-of-range + 개선조치 + 폐기·재작업·특채 → 3축 다값 & 출력 스키마 변경 — 재학습 아닌 재설계 & 처분 축만 제품화(C-02) · 3축 스키마 유상(C-01, 10 mm) · 무인·이물 판정 거절(C-30 커널 · C-22 미측정) \\ \hline
② 검사항목 식별자 & 도면 기반 DIM-\#\#\#\#, 마스터 217개 & 마스터가 없다 — 서식 기반 자유 문자열 & 매핑 대상이 존재하지 않는다 & 거절(C-05·06·07·36) — 마스터는 온담이 만든다. 상위 30개사 서식만 유상(C-08) · 분모 B 연계 유상(C-28) \\ \hline
③ 시간·수량 단위 & 로트 단위(초·중·종품 3점) & 연속 시계열 / 배치 / 건 & 1차 키가 다르다 & 건·계측기·오프라인 제품화(C-09·14·23) · 배치 앱 유상(C-40) · 시계열 거절(C-20 IoT) · 최소 연동 유상(C-18·21) \\ \hline
④ 규격 근거 & 도면 + 고객 사양 2중 & 법정 + 고객 + 자체 3중, 법정 우선 & 판정 우선순위 + 규제 리스크 & 자동 대조 제품화(C-10) · 법정 플래그 유상 + 판정 보조 한정 + ㊲ 배상 분리 연동(C-03) · 법정 DB 거절(C-11) \\ \hline
⑤ 부적합 처리 흐름 & 8D/NCR → 1차사 SQ 포털 & HACCP 개선조치 → 자체 보관 & Flow 존재 이유 절반 & 자체 보관 워크플로우·관청 리포트·폐기 코드 유상(C-04·17·27) · P1 발주 연동 거절(C-43) \\ \hline
\fi
\end{tabularx}}
\landscapeoff

%% ------------------------------------------------------------------ 3
\landscapeon
\secbar{3. 요구 47건 3분류 (C-01~47 · 별지 반복)}
\guide{항목·저자·차원(①~⑤)·공수(mm)·마진·유지보수 부담·분류(제품화/유상 커스텀/거절)·근거. 분류 기준 = 자동차 58개사에게도 가치가 있는가.}
{\footnotesize
\begin{tabularx}{\linewidth}{|L{14mm}|Y|L{20mm}|L{14mm}|L{18mm}|L{22mm}|L{22mm}|Y|}\hline
\hd{C-ID} & \hd{요구} & \hd{저자} & \hd{차원} & \hd{공수(mm)} & \hd{유지보수 부담} & \hd{분류} & \hd{근거} \\ \hline
\ifwbblank
\rd{8mm} &  &  &  &  &  &  &  \\ \hline
\rd{8mm} &  &  &  &  &  &  &  \\ \hline
\rd{8mm} &  &  &  &  &  &  &  \\ \hline
\rd{8mm} &  &  &  &  &  &  &  \\ \hline
\rd{8mm} &  &  &  &  &  &  &  \\ \hline
\rd{8mm} &  &  &  &  &  &  &  \\ \hline
\rd{8mm} &  &  &  &  &  &  &  \\ \hline
\rd{8mm} &  &  &  &  &  &  &  \\ \hline
\rd{8mm} &  &  &  &  &  &  &  \\ \hline
\rd{8mm} &  &  &  &  &  &  &  \\ \hline
\rd{8mm} &  &  &  &  &  &  &  \\ \hline
\rd{8mm} &  &  &  &  &  &  &  \\ \hline
\else
C-01 & 판정 라벨 3축 출력 스키마 & 나영선 & ① & 10 & H & 유상 & 스키마 재설계 — 별도 모듈. 처분 축은 C-02 \\ \hline
C-02 & 처분 코드 필드 & 나영선 & ① & 4 & M & 제품화 & 자동차 NCR 특채·재작업과 같다 \\ \hline
C-03 & CCP 한계 초과 자동 플래그 & 나영선 & ④ & 3 & M & 유상 & 판정 보조 한정 · ㊲ 항목 5 연동 \\ \hline
C-05 & 검사항목 마스터 구축 & 나영선 & ② & 7 & H & 거절 & 마스터는 고객의 데이터 거버넌스 \\ \hline
C-09 & 입고 검수 건 단위 기록 & 나영선 & ③ & 3 & M & 제품화 & 자동차 IQC도 건 단위 \\ \hline
C-20 & 냉장 온도 로거 시계열 수집·판정 & 한동석 & ③ & 8 & H & 거절 & IoT — GAUGE 범위 밖 \\ \hline
C-22 & 식품 이물 사진 판정 & 한동석 & ① & 6 & H & 거절 & 온담 세트 미측정 — ㊵ 항목 4 후 \\ \hline
C-27 & 폐기 사유 '온도 이탈' 코드 → RPT-INV-05 & 한동석 & ⑤ & 3 & M & 유상 & B-04 output — ㊵ \\ \hline
C-31 & 온프렘 배포 패키지 & 박세연 & — & 6 & H & 유상 & CR-044 보류분 → 온담 SOW \\ \hline
C-39 & 계정 1,200개 일괄 관리 & 박세연 & — & 3 & M & 제품화 & 김도윤의 '340개 협력사 한 화면' \\ \hline
C-41 & 점포 순위·경고 대시보드 & 오정근 & — & 2 & L & 거절 & 감시 도구화 — 커널 6-2 \\ \hline
C-47 & 사고 RCA 증빙 연 2회(제218문항) & 최영주 & — & 2 & L & 제품화 & ㉞ 항목 8 — 전 고객 공개 리포트 \\ \hline
\fi
\end{tabularx}}
\landscapeoff

%% ------------------------------------------------------------------ 4
\secbar{4. 분류별 계약 문구}
\guide{제품화(로드맵 편입·무상·시점 미약속) / 유상 커스텀(별도 SOW·단가·소유권) / 거절(대안 제시·사유). 각 분류에 한 문단.}
{\footnotesize
\begin{tabularx}{\linewidth}{|L{30mm}|Y|L{50mm}|}\hline
\hd{분류} & \hd{계약 문구} & \hd{리스크} \\ \hline
\ifwbblank
\rd{14mm}제품화 &  &  \\ \hline
\rd{14mm}유상 커스텀 &  &  \\ \hline
\rd{14mm}거절 &  &  \\ \hline
\else
제품화(14건 27 mm) & '표준 제품 로드맵에 편입, 라이선스 범위에서 무상 제공. 릴리스 시점은 분기 리뷰에서 통지하며 본 계약은 시점을 약속하지 않는다.' & 온담이 Now로 기억(태산 62항목) → 분기 리뷰 의사록에 미약속 재확인 \\ \hline
유상 커스텀(19건 59 mm) & '별도 SOW(59 mm × 1,200만 = 7.08억). 코드·문서 소유권 온담, 모델 가중치 딥게이지, 딥게이지는 온담 식별 정보를 제외한 재사용권을 가진다. 유지보수 연 SOW의 18\%.' & 온담 '21억 포함' 주장 → 공문 문구로 협상. SOW 없이 착수 = 138 mm 무상 \\ \hline
거절(14건 52 mm) & '마스터·법정 기준 DB·IoT 시계열·P1 포털·DR은 온담 IT본부 소관, 딥게이지는 인터페이스 사양 제공. 무인 판정·감시 대시보드는 설계 원칙과 충돌. 식품 이물 판정은 온담 세트 측정 후 재논의.' & '리소스 부족'으로 읽히면 '돈을 더 주면' → 사유는 소유·원칙·측정으로만 \\ \hline
\fi
\end{tabularx}}

%% ------------------------------------------------------------------ 5
\secbar{5. 용량 잠식 계산}
\guide{138 mm ÷ 개발 14명 = 9.9개월 = 연간 용량의 82.1\%. 커스텀만 받았을 때 / 제품화만 했을 때 / 거절 후의 세 시나리오.}
{\footnotesize
\begin{tabularx}{\linewidth}{|L{40mm}|L{22mm}|L{22mm}|Y|Y|}\hline
\hd{시나리오} & \hd{공수(mm)} & \hd{개발 개월} & \hd{자동차 58사 릴리스 공백} & \hd{판단} \\ \hline
\ifwbblank
\rd{10mm}47건 전부 수용 &  &  &  &  \\ \hline
\rd{10mm}분류 후(제품화 + 커스텀) &  &  &  &  \\ \hline
\rd{10mm}커스텀만 &  &  &  &  \\ \hline
\rd{10mm}거절 후 &  &  &  &  \\ \hline
\else
47건 전부 수용 & 138 & 9.9 & 분기 릴리스 3회 소멸 & 회사가 온담 SI가 된다 — 82.1\% \\ \hline
분류 후(제품화 + 커스텀) & 86 & 6.1 & 2회 소멸 — 상한 없으면 & 51.2\% → 상한 6 FTE(연 72 mm)를 계약 조건으로 \\ \hline
커스텀만 & 59 & 4.2 & 1회 & 35.1\% — 제품화 27은 어차피 한다(하한) \\ \hline
거절 후 & 0 & 0 & 없음 & 대한기공 2.4 · 코너스톤 소멸 · 판정자·측정 문제 없음 \\ \hline
\fi
\end{tabularx}}

%% ------------------------------------------------------------------ 6
\secbar{6. 집중 리스크 3표기}
\guide{ARR 비중 IR 기재 37.0\% / 검산 24.1\% / 개발 용량 82\%. 어느 것이 과대·과소인가, IR 덱에 어떻게 쓰는가.}
{\footnotesize
\begin{tabularx}{\linewidth}{|L{40mm}|L{20mm}|Y|L{22mm}|Y|}\hline
\hd{표기} & \hd{값} & \hd{분자/분모} & \hd{과대/과소} & \hd{IR 덱 문구} \\ \hline
\ifwbblank
\rd{10mm}ARR 비중(IR) &  &  &  &  \\ \hline
\rd{10mm}ARR 비중(검산) &  &  &  &  \\ \hline
\rd{10mm}개발 용량 잠식 &  &  &  &  \\ \hline
\else
ARR 비중(IR) & 37.0\% & 7.0 / 18.9(온담 연 청구 전액을 ARR로) & 과대 — 분자에 구축·운영 & '연 매출의 37\%' — 매출로는 쓰되 ARR로는 쓰지 않는다 \\ \hline
ARR 비중(검산) & 24.1\% & 3.78 / 15.68(라이선스만) & 적정(순SaaS 분모면 26.8\%) & 'ARR 15.68억 중 온담 3.78억(24.1\%)' \\ \hline
개발 용량 잠식 & 82.1\% & 138 / 168(14명 × 12) & 과소 — IR 덱에 없다 & '47건 138 mm — 분류 후 86 mm(51\%), 상한 6 FTE, 별도 SOW·로드맵 편입' \\ \hline
\fi
\end{tabularx}}

%% ------------------------------------------------------------------ 7
\secbar{7. BATNA 비교 — 대한기공}
\guide{연 2.4억 × 3년 / 커스텀 0 / 클라우드 / 4주. 매출·공수·도메인·텀시트 연동을 한 표에.}
{\footnotesize
\begin{tabularx}{\linewidth}{|L{34mm}|Y|Y|}\hline
\hd{기준} & \hd{온담 딜} & \hd{대한기공} \\ \hline
\ifwbblank
\rd{9mm}연 매출 / 3년 총액 &  &  \\ \hline
\rd{9mm}실질 ARR 증분 &  &  \\ \hline
\rd{9mm}공수(mm)·용량 &  &  \\ \hline
\rd{9mm}도메인(초기 SAM 안/밖) &  &  \\ \hline
\rd{9mm}클로징 기간·조건 &  &  \\ \hline
\rd{9mm}텀시트 연동(코너스톤) &  &  \\ \hline
\rd{9mm}리스크 &  &  \\ \hline
\else
연 매출 / 3년 총액 & 7.0 / 21.0억 & 2.4 / 7.2억(34.3\%) \\ \hline
실질 ARR 증분 & 3.78억(라이선스) & 2.40억(전액 SaaS) \\ \hline
공수(mm)·용량 & 138(분류 후 86) · 82.1\%(51.2\%) & 0 \\ \hline
도메인(초기 SAM 안/밖) & 확장 SAM(식품 — 제조사 아님) & 초기 SAM 안(1차 협력사) — 제품이 흔들리지 않는다 \\ \hline
클로징 기간·조건 & 47건 협상 + CFO 3주 → D+900 전후 · 2법인 & 표준계약 4주 · 클라우드 \\ \hline
텀시트 연동(코너스톤) & 선행 조건 충족 & 온담 거절 = 코너스톤 소멸 — 한강 단독 \\ \hline
리스크 & 용량·법정 기준 오판정·배상 100\%·판정자 미지정·미측정 & ACV 낮음 · 채널 의존(김도윤) \\ \hline
\fi
\end{tabularx}}

%% ------------------------------------------------------------------ 8
\secbar{8. Go / No-go 권고와 조건}
\guide{받는다면 어떤 조건에서, 거절한다면 무엇을 대신. '자동차 버티컬을 벗어날 것인가'에 대한 답을 한 문단.}
\begin{fieldbox}
\ifwbblank
\lines{7}{9.5mm}
\else
조건부 Go — 일곱 조건이 계약 문장으로 들어갈 때만: ① 47건 3분류(제품화 14 / 유상 19 / 거절 14)를 계약 별지로 ② 딥게이지 투입 상한 6 FTE(연 72 mm) ③ 유상 커스텀 별도 SOW 7.08억 ④ 데이터 이관 정본 판정자 나영선 지정, 커버리지 38.6\% 미매핑분은 온담 마스터 정비 후 ⑤ 온담 4세트 80건 측정 계획을 인수 기준 조항으로(㊵ 항목 4) ⑥ 배상 상한 4구간(㊲ 항목 5), AI 판정·법정 기준 판정 제외 ⑦ 61점 라인 정의는 계약 별지. 그리고 대한기공 표준계약을 병행 체결한다 — BATNA를 실제로 만든다. 자동차를 벗어날 것인가: 벗어나지 않는다. 온담은 '자동차 제품이 식품에서도 서는가'를 측정하는 한 건이고, 두 번째 식품 고객은 측정 결과 이후에 판단한다. 조건 ①②⑤ 중 하나라도 빠지면 No-go — 대한기공 + 한강.
\fi
\end{fieldbox}

\end{document}